\chapter{Open problems}\label{ch:open}

Each problem below is stated precisely enough to be refuted or proved, with the evidence currently available and the obstruction to a proof as far as we know it. None is a claim. Problems are grouped by the part of the monograph from which they arise. The label ``computational evidence'' always refers to the stated finite corpus.

\section{Foundations and ancestry}
\begin{problem}[growth of ancestry]\label{prob:growth}
For the content $72^r$ (the $r$-fold repetition of the labels of the worked example), does $A(72^r)/r$ converge as $r\to\infty$, and is $|F_{72^r}|$ asymptotically exponential in $r$? \emph{Evidence.} Exact values for $r\le7$ (re-run) and $r=8,\dots,10$ (inherited, \texttt{NOT\_RERUN}); the elementary lower bound is far from sharp (\cref{ch:ancestry}). \emph{Obstruction.} The set of histories mixes reorderings from the source $1$ with histories through other sources; no transfer-matrix description of the recursion is known.
\end{problem}
\begin{problem}[three notions of ambiguity]\label{prob:source}
Find sharp relations among $A(t)$, $A(t\,|\,m_0)$ and $\Ent(M|T)$ beyond \cref{thm:hierarchy,prop:entropy-bounds}. \emph{Evidence.} The three can differ by more than a constant factor on the worked example. \emph{Obstruction.} The weights of different sources are not controlled by the local rule.
\end{problem}
\begin{problem}[evolution of collective fields]\label{prob:evolution}
Extend the framework from one history of one finite set of events to the generative evolution of collective fields over many interacting sets. \emph{Status.} Not attempted in this monograph; no definition of the object is proposed here. We list it only because \cref{ch:intro} disclaims it.
\end{problem}

\section{Arithmetic topology and the atlas}
\begin{problem}[closure of the atlas]\label{prob:atlas-closure}
Is the atlas of $32$ survival-law classes of \cref{ch:atlas} complete for the cycles it describes, or do further classes appear on enlarging the search? \emph{Evidence.} The atlas is a bounded search; the held-out prediction record of \cref{obs:atlas-pred} supports the classification on the held-out data but cannot prove closure. \emph{Obstruction.} The families were computed by dualisation over a finite pool of sources, and no a priori bound on the number of classes is known.
\end{problem}
\begin{problem}[representative dependence]\label{prob:representatives}
Classes in the atlas are defined by low-weight representatives of homology classes. Is the survival law an invariant of the homology class, i.e.\ do all representatives of a class give the same law? \emph{Evidence.} Observed on the sampled representatives only. \emph{Note.} The law for a fixed cycle is proved (\cref{thm:AT1}); the question is about the class.
\end{problem}
\begin{problem}[realisability]\label{prob:realisability}
Which of the survival laws realised on $P_6$ are realised by contents with prescribed prime support, and which abstract lattice-valued laws are never realised? \emph{Evidence.} CRT realisation in the tested range (\cref{prop:CRT}). \emph{Obstruction.} The constraint is a Diophantine condition on the labels.
\end{problem}

\section{Dimension raising}
\begin{problem}[non-degeneracy of off-diagonal lifts]\label{prob:nondeg}
Prove that every off-diagonal lift $L_{c_1,c_2}(z)$ of a non-trivial class is non-trivial (the statement ND-1 of the source programme). \emph{Evidence.} Finite-exact on the censused lifts of the anchor (\cref{obs:nondeg}). \emph{Obstruction.} The boundary \cref{lem:LT0} gives the cycle property but not injectivity of the lift on homology.
\end{problem}
\begin{problem}[every $H_3$ class is a lift?]\label{prob:every-lift}
For a seven-label content, is every arithmetic $H_3$ class a lift of a six-label class? \emph{Evidence.} Only lifted patterns were censused (\cref{prop:H3-examples}); no non-lift class is known and none has been excluded.
\end{problem}
\begin{problem}[injectivity of $\first_\ast$]\label{prob:Jinjective}
Is the map induced by $\first$ on homology injective on the classes in \cref{thm:JX}? \emph{Evidence.} Facet and retraction criteria (J-1, J-2) are proved; injectivity is verified on the finite cases (J-3--J-5) only.
\end{problem}
\begin{problem}[propagation of filling dependence]\label{prob:propagation}
Does the filling dependence of \cref{thm:filling-dependence} propagate to every dimension under iterated lifts, so that one gets a Dimension-Raising Obstruction Theorem? \emph{Evidence.} Proved for one lift (HL-1) and verified for two (HL-2); the three-lift search (HL-3) is inherited and \texttt{NOT\_RERUN}.
\end{problem}
\begin{problem}[role-invariant lifting]\label{prob:rolemap}
Is the statement that a lift does not require the lower-dimensional hole to survive true for every role map, and not only for the identity role map used in the seven-label census? \emph{Evidence.} Supported for the identity role map and the two sub-contents examined (\cref{obs:rolemap}); the role-invariant version is open.
\end{problem}

\section{Information recovery}
\begin{problem}[does radius three suffice?]\label{prob:radius3}
Is the fine ball signature through radius three a complete invariant of the Betti vector for distinct-label pieces for every $k\le5$ and all label ranges, and does it fail or hold at $k=6$? \emph{Evidence.} No counterexample on $k\le5$ and labels $\le26$ in the arithmetic census; radius two provably fails (witnesses W5, W6). This is the first information-recovery conjecture of the programme and is a conjecture, not a result.
\end{problem}
\begin{problem}[arithmetic origin of the $(Q,S_{s2})$ separation]\label{prob:arith-origin}
The separation of positive and negative systems by the pair $(Q,S_{s2})$ holds on every arithmetic regime tested but fails as a general principle (\cref{obs:S2-falsified}, prefix-family counterexample). Does arithmetic closure force it, and if so by what mechanism? \emph{Obstruction.} The counterexamples are prefix-defined but not arithmetic, so any proof must use divisibility.
\end{problem}

\section{Forbidden histories and locality}
\begin{problem}[local criterion for hollowness]\label{prob:local-hollow}
Characterise, from the dead-set nerve and local data alone, which positive systems have $\beta_2>0$. \emph{Evidence.} Counts for the corpus and the observation that $S_{s2}$ is smaller for the positive system in $4/4$ equal-signature classes (\cref{obs:cycle-census}); no criterion is proved.
\end{problem}
\begin{problem}[constraint locality beyond six labels]\label{prob:locality-k}
Does the pairwise-adjacency theorem (\cref{thm:pairwise}) hold for $k\ne6$? If so, is there a proof that does not reduce to a finite verification for each $k$? \emph{Evidence.} None beyond $k=6$: the theorem is proved for six labels by a computer-assisted reduction, and we make no claim for other $k$. For the triple constraint the analogous statement is false at $k=6$ (\cref{prop:triple}).
\end{problem}

\section{Framework}
\begin{problem}[relation to cellular resolutions]\label{prob:resolution}
Is there a relation between the survival law \cref{thm:AT1} and the minimal free resolution of the monomial module generated by the vertex labels (\cref{sec:lcm})? In particular, does $\beta_2(X_C(m))>0$ force a non-zero Betti number of the module in multidegree $m$? \emph{Status.} Not examined: no resolution has been computed.
\end{problem}
